\notechapter{N-20230117}{More Function Exercises: Periodicity, Piecewise Definitions, Images, and Functional Equations}
\enlargethispage{4\baselineskip}

\section{Continuing the function chapter}

This note continues the function exercises from the previous day.  The pages mix handwritten solutions with printed exercises.  The recurring theme is to translate each problem into one of several languages:
\[
  \text{formula},\qquad \text{graph},\qquad \text{image/preimage},\qquad \text{iteration},\qquad \text{functional equation}.
\]
Many errors in such problems come from using a formula while ignoring its domain or from treating a functional equation as if it held only for special values.

\section{Piecewise functions}

The first page includes piecewise-defined functions.  For example, a function may be defined differently on intervals such as
\[
  x<0,\qquad x=0,\qquad x>0.
\]
When studying continuity or solving equations involving such a function, one must first decide which case the input belongs to.  If the input itself is an expression such as $f(x)$ or $x+1$, the relevant case may change during the calculation.

The safe method is:
\begin{enumerate}[(i)]
  \item identify the possible regions of the input;
  \item solve the equation in each region;
  \item check that the solution actually belongs to the assumed region.

\end{enumerate}

\section{Periodicity from functional equations}

A relation such as
\[
  f(x+1)=1-f(x)
\]
implies
\[
  f(x+2)=1-f(x+1)=f(x).
\]
Thus the function is periodic with period $2$.  The note uses this style of argument in several places: apply the same equation repeatedly until the input returns to a familiar form.

If instead
\[
  f(x+T)=f(x)
\]
for all $x$, then $T$ is a period.  The smallest positive period, if it exists, is called the fundamental period.

\section{Graph transformations}

The pages include small sketches of graphs.  The standard transformations are:
\[
  y=f(x-a) \quad \text{shift right by }a,
\]
\[
  y=f(x)+b \quad \text{shift up by }b,
\]
\[
  y=-f(x) \quad \text{reflect across the }x\text{-axis},
\]
\[
  y=f(-x) \quad \text{reflect across the }y\text{-axis}.
\]
Absolute values fold graphs:
\[
  y=|f(x)|
\]
reflects the negative part upward, while
\[
  y=f(|x|)
\]
copies the right half to the left.

\section{Image of an interval}

For a continuous function on an interval, the image is also an interval.  This is the intermediate value principle in set language.  If the function is monotone on the interval, then the endpoints determine the image directly.

The note's exercises often ask for the range of an expression.  The method is to rewrite it as a function of one variable, find monotonicity or extrema, and then translate the result back into an interval.

\section{Functional equations and substitution}

The handwritten solutions repeatedly substitute special values such as $x=0$, $y=0$, or $y=x$.  This is a standard strategy.  For a functional equation in two variables,
\[
  F(x,y,f(x),f(y),f(x+y))=0,
\]
choosing simple values isolates pieces of the function.

For example, if
\[
  f(x+y)=f(x)+f(y)+c,
\]
then setting $x=y=0$ determines $f(0)$, and subtracting a constant can reduce the equation to the additive Cauchy equation.

\section{Counting functions between finite sets}

The printed exercise page also touches finite functions.  If $|A|=m$ and $|B|=n$, then the number of functions $A\to B$ is
\[
  n^m,
\]
because each element of $A$ independently chooses one image in $B$.

If one requires injectivity and $m\le n$, the number is
\[
  n(n-1)\cdots(n-m+1).
\]
If one requires bijectivity and $m=n$, the number is
\[
  n!.
\]

\section{A compact bridge}

The note is a training page for reading functions flexibly.  A function may be a formula, a graph, a rule on a finite set, or an unknown object constrained by an equation.  Each representation suggests a different method: casework for piecewise definitions, transformations for graphs, counting for finite domains, and substitution for functional equations.



\paragraph{Expanded synthesis.}
For this chapter, the elementary material should be read as a study of what remains invariant when coordinates change. The earlier sections (`Continuing the function chapter`, `Piecewise functions`, `Periodicity from functional equations`, `Graph transformations`, `Image of an interval`) compute with arrays, bases, pivots, determinants, or eigenvectors; the deeper object is a linear map together with the spaces it acts on. A matrix is a representation of that map after bases have been chosen. This is why formulas such as
\[
  A' = P^{-1}AP,\qquad \widetilde A = T^{-1}AS
\]
are not cosmetic: they distinguish an intrinsic morphism from a coordinate description.

The structural hierarchy is as follows. Kernels record what a map forgets, images record what it reaches, cokernels record obstructions to solvability, and quotients record information after an equivalence has been imposed. Determinants act on the top exterior power and therefore measure oriented volume; traces are invariant under cyclic rearrangement and become characters in representation theory; adjoints depend on an inner product and lead to self-adjoint spectral theory.

A useful way to return to the computations is to ask, for every formula, which invariant it is measuring. Row reduction measures rank and kernel; diagonalization measures decomposition into invariant subspaces; Gram--Schmidt measures orthogonal projection; positive definiteness measures ellipsoid geometry; tensor notation measures how components transform. The elementary methods are therefore the finite-dimensional laboratory in which duality, quotienting, functoriality, and spectral decomposition can be seen clearly.


\section{Functional equations as dynamics}

A functional equation often specifies a transformation rule rather than a direct formula.  For example, equations involving $f(x+1)$, $f(2x)$, or $f(f(x))$ are naturally studied by iteration.  The orbit of $x$ under a transformation $T$ is
\[
  x,\ T(x),\ T^2(x),\ \ldots.
\]
A functional equation may prescribe how $f$ behaves along these orbits.

For instance, if
\[
  f(x+1)=f(x),
\]
then $f$ is constant along orbits of the translation $T(x)=x+1$.  Periodicity is invariance under a group action of $\mathbb Z$.

\section{Cauchy's equation and regularity}

Cauchy's equation is
\[
  f(x+y)=f(x)+f(y).
\]
If $f$ is continuous, measurable, or bounded on an interval, then
\[
  f(x)=cx.
\]
But without regularity assumptions, there are wild solutions constructed using a Hamel basis of $\mathbb R$ over $\mathbb Q$.

This shows why elementary functional equations often include hidden conditions such as continuity or monotonicity.  These conditions are not technical decorations; they exclude pathological solutions arising from the axiom of choice.

\section{Symbolic dynamics}

Piecewise-defined functions and iteration can lead to symbolic dynamics.  If an interval map sends subintervals across one another, one can encode an orbit by a sequence of symbols recording which subinterval contains each iterate.

The shift map
\[
  \sigma((a_0,a_1,a_2,\ldots))=(a_1,a_2,a_3,\ldots)
\]
is a central model.  Complicated interval dynamics can often be compared to shifts on symbol spaces.

Thus a piecewise function graph in an elementary note can be the beginning of chaos, coding of orbits, and topological entropy.

\section{Rigidity from regularity}

Many functional equations have many algebraic solutions but few regular solutions.  The pattern is:

\begin{center}
\small
\begin{tabularx}{\textwidth}{@{}lYY@{}}
\toprule
Equation & Algebraic freedom & Regularity conclusion \\
\midrule
$f(x+y)=f(x)+f(y)$ & Hamel-linear maps & $f(x)=cx$ under continuity \\
$f(xy)=f(x)+f(y)$ & Logarithmic freedom over the multiplicative group & $f(x)=c\log x$ under continuity \\
$f(x+T)=f(x)$ & Arbitrary values on a fundamental domain & Periodic structure under smoothness \\
\bottomrule
\end{tabularx}
\end{center}

The original piecewise and functional-equation exercises should therefore be read as local training for a general principle: algebraic equations become meaningful only after the allowed regularity class is specified.

% Second-pass deepening for waves 2-4: wave 3 A-C part 2
\section{Piecewise functions and partitions}

A piecewise function is a function defined after partitioning the domain into regions.  The mathematical issue is not only which formula applies, but how the pieces meet.  Discontinuities occur at boundaries of the partition; differentiability requires matching first-order behavior across boundaries.

For example, if
\[
  f(x)=\begin{cases}f_1(x),&x<a,\\ f_2(x),&x\ge a,
  \end{cases}
\]
then continuity at $a$ requires
\[
  \lim_{x\to a^-}f_1(x)=f_2(a).
\]
Differentiability additionally requires matching slopes.  This is a local gluing condition.

\section{Functional equations and invariants}

A functional equation often says that $f$ is invariant or equivariant under transformations.  The equation
\[
  f(x+T)=f(x)
\]
says $f$ is invariant under the translation action of $T\mathbb Z$.  The equation
\[
  f(ax)=bf(x)
\]
says $f$ is homogeneous with respect to scaling.

Thus solving a functional equation often means finding the invariants of a group or semigroup action.  This perspective makes the original examples more natural: periodicity, parity, and scaling are all symmetry constraints.
